\documentclass[12pt,oneside]{book}

% ============================================================
% METADATOS
% ============================================================
\newcommand{\universidad}{Universidad de Buenos Aires (UBA)}
\newcommand{\facultad}{Facultad de Derecho}
\newcommand{\materia}{Derecho Procesal Civil y Comercial}
\newcommand{\titulo}{Modelo de Demanda Civil: Daños y Perjuicios, Citación en Garantía, Prueba y Preclusión}
\newcommand{\autor}{Isaac Edgar Camacho Ocampo}
\newcommand{\anio}{2024}

% ============================================================
% PAQUETES
% ============================================================
\usepackage[utf8]{inputenc}
\usepackage[T1]{fontenc}
\usepackage[spanish,es-nodecimaldot]{babel}

\usepackage[
    top=2.5cm,
    bottom=2.5cm,
    left=3cm,
    right=2.5cm,
    headheight=15pt
]{geometry}

\usepackage{amsmath}
\usepackage{amssymb}
\usepackage{mathtools}

\usepackage{graphicx}
\usepackage{float}
\usepackage{caption}
\usepackage{subcaption}

\usepackage{booktabs}
\usepackage{array}
\usepackage{longtable}
\usepackage{tabularx}

\usepackage{enumitem}

\usepackage{xcolor}
\usepackage[most]{tcolorbox}

\usepackage{fancyhdr}
\usepackage{ragged2e}

\usepackage[
    colorlinks=true,
    linkcolor=blue!50!black,
    citecolor=green!40!black,
    urlcolor=blue!60!black,
    pdftitle={\titulo},
    pdfauthor={\autor},
    pdfsubject={Apuntes universitarios},
    bookmarks=true
]{hyperref}

% ============================================================
% CONFIGURACIÓN
% ============================================================
\graphicspath{
    {./img/}
    {./graficos/}
    {./figuras/}
}

\setcounter{tocdepth}{2}

\setlength{\parindent}{0pt}
\setlength{\parskip}{0.5em}

\setlist[itemize]{
    topsep=0.3em,
    itemsep=0.2em
}

\setlist[enumerate]{
    topsep=0.3em,
    itemsep=0.2em
}

\clubpenalty=10000
\widowpenalty=10000

% ============================================================
% COMANDOS MATEMÁTICOS
% ============================================================
\newcommand{\abs}[1]{\left|#1\right|}
\newcommand{\norm}[1]{\left\lVert#1\right\rVert}
\newcommand{\vect}[1]{\mathbf{#1}}
\newcommand{\deriv}[2]{\frac{d#1}{d#2}}
\newcommand{\pderiv}[2]{\frac{\partial#1}{\partial#2}}

% ============================================================
% ENTORNOS / CAJAS ACADÉMICAS
% ============================================================
\newtcolorbox{definition}[1]{
    enhanced,
    breakable,
    title={Definición: #1},
    fonttitle=\bfseries,
    colback=blue!3,
    colframe=blue!50!black,
    boxrule=0.8pt,
    arc=2mm
}

\newtcolorbox{important}{
    enhanced,
    breakable,
    title={Importante},
    fonttitle=\bfseries,
    colback=yellow!8,
    colframe=orange!70!black,
    boxrule=0.8pt,
    arc=2mm
}

\newtcolorbox{example}[1]{
    enhanced,
    breakable,
    title={Ejemplo: #1},
    fonttitle=\bfseries,
    colback=green!4,
    colframe=green!40!black,
    boxrule=0.8pt,
    arc=2mm
}

\newtcolorbox{formula}[1]{
    enhanced,
    breakable,
    title={Fórmula / Regla: #1},
    fonttitle=\bfseries,
    colback=purple!4,
    colframe=purple!50!black,
    boxrule=0.8pt,
    arc=2mm
}

\newtcolorbox{exercise}[1]{
    enhanced,
    breakable,
    title={Ejercicio: #1},
    fonttitle=\bfseries,
    colback=gray!5,
    colframe=gray!60!black,
    boxrule=0.8pt,
    arc=2mm
}

\newtcolorbox{note}{
    enhanced,
    breakable,
    title={Nota},
    fonttitle=\bfseries,
    colback=white,
    colframe=gray!50,
    boxrule=0.6pt,
    arc=2mm
}

% ============================================================
% CONFIGURACIÓN FINAL (ENCABEZADOS Y PIES)
% ============================================================
\pagestyle{fancy}
\fancyhf{}

\fancyhead[L]{\nouppercase{\leftmark}}
\fancyhead[R]{\materia}

\fancyfoot[C]{\thepage}

\renewcommand{\headrulewidth}{0.4pt}
\renewcommand{\footrulewidth}{0pt}

\fancypagestyle{plain}{
    \fancyhf{}
    \fancyfoot[C]{\thepage}
    \renewcommand{\headrulewidth}{0pt}
}

% ============================================================
% DOCUMENTO
% ============================================================
\begin{document}

% ------------------------------------------------------------
% PORTADA
% ------------------------------------------------------------
\begin{titlepage}

\thispagestyle{empty}

\begin{center}

\vspace*{1cm}

{\large \textsc{\universidad}}

\vspace{0.2cm}

{\large \textsc{\facultad}}

\vspace{3cm}

{\Huge \textbf{\materia}}

\vspace{1cm}

{\LARGE \titulo}

\vspace{0.8cm}

\textit{Comprender antes que memorizar es el camino más corto hacia el conocimiento verdadero.}

\vspace{1.2cm}

\rule{0.70\textwidth}{0.5pt}

\vspace{0.5cm}

{\large Apuntes de estudio}

\vspace{0.5cm}

\rule{0.70\textwidth}{0.5pt}

\vfill

{\large \autor}

\vspace{0.3cm}

{\large \anio}

\end{center}

\vfill

\begin{flushleft}
\textit{Este es un modesto aporte a los estudiantes de ``Derecho Procesal Civil y Comercial'' no pretende sustituir el material oficial de la materia.}
\end{flushleft}

\end{titlepage}

% ------------------------------------------------------------
% ÍNDICE
% ------------------------------------------------------------
\tableofcontents

% ------------------------------------------------------------
% CONTENIDO
% ------------------------------------------------------------

\chapter{Introducción general}

\section{Contexto y caso unificador de la clase}

La clase se organiza en torno a un único caso que atraviesa transversalmente todos los temas tratados: un kayakista fallece a raíz de la colisión de su embarcación con una lancha en el Río Tigre. A partir de este caso unificador, se recorre cada uno de los pasos que un abogado debe dar para llevar un hecho de esta naturaleza hasta la instancia de juicio: desde la redacción de la demanda hasta el ofrecimiento de la prueba.

El eje práctico de la clase es la exposición, por parte de un alumno, de un modelo de demanda por daños y perjuicios construido sobre ese caso, el cual es revisado y corregido en el desarrollo de la exposición conforme a los requisitos legales exigibles.

\section{Relevancia profesional de los temas tratados}

El conocimiento de la estructura correcta de la demanda, prevista en el \textbf{artículo 330 del Código Procesal Civil y Comercial de la Nación (CPCCN)}, resulta indispensable para cualquier abogado, con independencia de su especialidad, dado que todo profesional del derecho redacta o contesta demandas en algún momento de su ejercicio. Un escrito con defectos formales puede ser rechazado o utilizado como defensa por la contraparte.

La \textbf{citación en garantía} prevista en el artículo 118 de la Ley de Seguros es una herramienta de uso cotidiano en casos de daños de tránsito, accidentes laborales y, en general, en todo hecho cubierto por un seguro. Desconocerla puede implicar dejar sin cobertura al cliente o perder la posibilidad de accionar contra el asegurador.

Distinguir el \textbf{daño moral} del \textbf{daño psicológico} no constituye una cuestión meramente teórica: de esa distinción depende qué prueba se produce, qué perito se ofrece y, en definitiva, el monto que percibirá el cliente. Un error en este punto reduce directamente la indemnización.

La \textbf{preclusión procesal} es la disciplina que ordena el proceso. Comprender que un acto procesal cierra una etapa —incluso si el plazo legal aún no venció— evita errores irreversibles, como omitir el ofrecimiento de testigos o de prueba documental crucial.

Finalmente, los \textbf{rubros indemnizatorios} (lucro cesante, pérdida de chance, daño emergente, daño psicológico, gastos médicos, entre otros) constituyen el eje económico de todo juicio civil. Identificarlos correctamente y acreditarlos con la prueba adecuada define, en gran medida, el resultado económico del proceso para el cliente.

\section{Temas tratados en la clase}

\begin{itemize}
    \item Estructura y redacción de una demanda civil (Art. 330 CPCCN).
    \item Objeto, legitimación activa y pasiva en la demanda.
    \item Citación en garantía del asegurador (Art. 118 Ley de Seguros).
    \item Competencia civil vs. penal: tribunales nacionales y federales.
    \item Rubros indemnizatorios: daño patrimonial y extrapatrimonial.
    \item Daño moral vs. daño psicológico: diferencias y prueba.
    \item Medios de prueba: documental, testimonial, informativa, pericial.
    \item Preclusión procesal y hecho nuevo.
    \item Cargas procesales vs. obligaciones.
\end{itemize}

\section{Mapa temático de la clase}

El siguiente cuadro sintetiza los bloques temáticos desarrollados en la clase, agrupando cada tema con sus palabras clave y la pregunta central que permite repasarlo.

\begin{longtable}{>{\raggedright\arraybackslash}p{0.06\textwidth} >{\raggedright\arraybackslash}p{0.22\textwidth} >{\raggedright\arraybackslash}p{0.26\textwidth} >{\raggedright\arraybackslash}p{0.34\textwidth}}
\toprule
\textbf{N.\textsuperscript{o}} & \textbf{Tema} & \textbf{Palabras clave} & \textbf{Pregunta clave} \\
\midrule
\endfirsthead
\toprule
\textbf{N.\textsuperscript{o}} & \textbf{Tema} & \textbf{Palabras clave} & \textbf{Pregunta clave} \\
\midrule
\endhead
\multicolumn{4}{l}{\textit{Bloque I: La demanda civil (temas 1--4)}} \\
\midrule
1 & Estructura de la demanda & Art. 330 CPCCN --- requisitos formales & ¿Qué debe contener obligatoriamente una demanda? \\
2 & Objeto de la demanda & Reclamación --- monto --- exactitud & ¿Qué es el objeto y por qué debe ser preciso? \\
3 & Legitimación activa y pasiva & Actor --- demandado --- calidad jurídica & ¿Quién puede demandar y a quién se puede demandar? \\
4 & Competencia & Nacional --- federal --- penal --- civil & ¿Ante qué juez presento la demanda? \\
\midrule
\multicolumn{4}{l}{\textit{Bloque II: Seguro y citación en garantía (temas 5--6)}} \\
\midrule
5 & Citación en garantía & Art. 118 --- seguro obligatorio --- parte procesal & ¿Por qué cito al asegurador y cómo lo hago? \\
6 & Contrato de seguro & Cobertura --- riesgo --- siniestro --- póliza & ¿Cuándo responde el seguro y por qué debe ser parte? \\
\midrule
\multicolumn{4}{l}{\textit{Bloque III: Rubros indemnizatorios (temas 7--10)}} \\
\midrule
7 & Daño patrimonial & Lucro cesante --- pérdida de chance --- emergente & ¿Qué pérdidas económicas puedo reclamar? \\
8 & Daño extrapatrimonial & Daño moral --- espiritual --- inmaterial & ¿Cómo se cuantifica el sufrimiento? \\
9 & Daño psicológico & Perito --- terapia --- cuantificación médica & ¿En qué se diferencia del daño moral? \\
10 & Gastos y otros rubros & Hospital --- velatorio --- sepelio --- obra social & ¿Qué gastos concretos son reclamables? \\
\midrule
\multicolumn{4}{l}{\textit{Bloque IV: Prueba procesal (temas 11--16)}} \\
\midrule
11 & Prueba documental & Soporte --- acta --- instrumento --- video & ¿Qué documentos debo acompañar con la demanda? \\
12 & Prueba testimonial y pericial & Testigos --- expertos --- preguntas --- dictamen & ¿Cómo se producen estas pruebas? \\
13 & Prueba informativa & Oficio --- organismos --- dependencias públicas & ¿Cuándo pido información a terceros? \\
14 & Hecho nuevo & Requisitos --- plazo --- audiencia 360 & ¿Qué sucede si aparece prueba nueva después de la demanda? \\
15 & Cargas procesales & Sin sanción --- oportunidad --- estrategia & ¿Qué diferencia una carga de una obligación? \\
16 & Preclusión procesal & Cierre de etapa --- vencimiento --- acto cumplido & ¿Qué pasa si contesté antes de que venza el plazo? \\
\bottomrule
\end{longtable}

\chapter{La demanda civil: estructura, objeto y sujetos (Art. 330 CPCCN)}

\section{Objeto de la demanda}

El objeto de la demanda consiste en el reclamo de daños y perjuicios dirigido contra los siguientes sujetos: el patrón de la lancha, el marinero de la lancha, la empresa de transporte marítimo (empresa de colectivos fluviales) y la compañía de seguros, esta última citada en garantía en virtud del artículo 118 de la Ley de Seguros.

\begin{example}{Redacción del objeto de la demanda}
``Vengo a promover demanda por daños y perjuicios contra los señores patrón y marinero, y también contra la empresa de colectivos, por ser responsables de ocasionar lesiones graves que causaron la muerte de quien fuera mi mujer y compañía de trabajo, la señora [nombre], y cito en garantía a la compañía de seguros en virtud del artículo 118 de la Ley de Seguros.''
\end{example}

% TODO: el apunte original no consigna el nombre de la víctima, reemplazado por el propio expositor con ``[nombre]''.

\begin{important}
No es indispensable indicar el monto reclamado dentro del objeto de la demanda; ello puede precisarse en otro momento procesal. Sin embargo, el monto o, al menos, una pauta para estimarlo no puede faltar del escrito de demanda, ya que su ausencia total habilitaría a la contraparte a interponer una excepción de defecto legal. Precisar claramente qué se reclama permite además que la parte demandada evalúe su estrategia de defensa, incluyendo la posibilidad de allanarse.
\end{important}

\section{Requisitos formales: Art. 330 CPCCN}

El artículo 330 del Código Procesal Civil y Comercial de la Nación establece los requisitos formales que debe reunir toda demanda.

\begin{formula}{Art. 330 CPCCN --- Contenido de la demanda}
La demanda será deducida por escrito y contendrá:
\begin{enumerate}[label=\arabic*.\textsuperscript{o})]
    \item El nombre y domicilio del demandante.
    \item El nombre y domicilio del demandado.
    \item La cosa demandada, designándola con toda exactitud.
    \item Los hechos en que se funde, explicados claramente.
    \item El derecho expuesto sucintamente, evitando repeticiones innecesarias.
    \item La petición en términos claros y positivos. La demanda deberá precisar el monto reclamado, salvo cuando al actor no le fuera posible determinarlo al promoverla por las circunstancias del caso o porque la estimación dependiera de elementos aún no definitivamente fijados y la promoción de la demanda fuese imprescindible para evitar la prescripción de la acción.
\end{enumerate}
\end{formula}

\subsection{Verificación del cumplimiento de cada inciso}

A partir del modelo de demanda expuesto, se verificó el cumplimiento de cada uno de los incisos del artículo 330:

\begin{longtable}{>{\raggedright\arraybackslash}p{0.42\textwidth} >{\raggedright\arraybackslash}p{0.48\textwidth}}
\toprule
\textbf{Inciso del Art. 330} & \textbf{¿Cumplido por el expositor?} \\
\midrule
\endfirsthead
\toprule
\textbf{Inciso del Art. 330} & \textbf{¿Cumplido por el expositor?} \\
\midrule
\endhead
Inc. 1.\textsuperscript{o} --- Nombre y domicilio del demandante & Sí: presenta a la señora Pérez como cónyuge supérstite. \\
Inc. 2.\textsuperscript{o} --- Nombre y domicilio del demandado & Sí: patrón, marinero, empresa fluvial y aseguradora. \\
Inc. 3.\textsuperscript{o} --- La cosa demandada con exactitud & Sí: reclamación por daños y perjuicios con detalle. \\
Inc. 4.\textsuperscript{o} --- Los hechos en que se funde & Sí: relata el accidente con la lancha en el Tigre. \\
Inc. 5.\textsuperscript{o} --- El derecho (sucintamente) & A desarrollar con las normas aplicables. \\
Inc. 6.\textsuperscript{o} --- Petición clara y monto & Puede diferirse si hay imposibilidad, con fundamento. \\
\bottomrule
\end{longtable}

\section{Legitimación activa y pasiva}

\subsection{Cargas procesales vs. obligaciones}

En el derecho procesal no existen, en sentido estricto, ``obligaciones'', sino \textbf{cargas}.

\begin{definition}{Carga procesal}
La carga procesal no lleva aparejada una sanción. La obligación, en cambio, es el deber de hacer o no hacer, exigible judicialmente y susceptible de ser compelida o forzada por vía coactiva. La carga, si no se ejerce, no genera una sanción sino la pérdida de una oportunidad: si no se ofrecen testigos, lo que se pierde es la oportunidad de valerse de esa prueba. En el proceso civil, todos los actos que las partes deben realizar son cargas: en la medida en que se ejerzan, se incrementan las probabilidades de obtener un resultado favorable.
\end{definition}

\subsection{Legitimación activa}

La legitimación activa corresponde, en el caso analizado, al cónyuge supérstite, en tanto era el esposo de la persona fallecida. Plantear expresamente la legitimación activa desde el inicio de la demanda constituye una estrategia procesal que anticipa y neutraliza una eventual excepción previa de falta de legitimación activa que pudiera oponer la parte demandada. Explicitar con claridad el fundamento de la legitimación equivale a adelantarse a la defensa de la contraparte.

\subsection{Legitimación pasiva: distintos títulos de responsabilidad}

Todos los sujetos demandados resultan responsables por el mismo hecho, aunque no por el mismo título jurídico. El marinero y el patrón responden como responsables directos por su negligencia; la empresa de transporte responde por ser dueña de la cosa riesgosa (la lancha), entendiéndose por cosa riesgosa cualquier objeto de esa naturaleza, como un automóvil o una embarcación; y la compañía de seguros responde por el contrato que la vincula con la relación jurídica, en virtud del artículo 118 de la Ley de Seguros.

\begin{longtable}{>{\raggedright\arraybackslash}p{0.28\textwidth} >{\raggedright\arraybackslash}p{0.62\textwidth}}
\toprule
\textbf{Demandado} & \textbf{Título de responsabilidad} \\
\midrule
\endfirsthead
\toprule
\textbf{Demandado} & \textbf{Título de responsabilidad} \\
\midrule
\endhead
Marinero & Responsable directo por negligencia: no estaba en la popa verificando las condiciones para la maniobra de zarpada. \\
Patrón & Responsable directo por negligencia: dio la orden o permitió la maniobra en condiciones inseguras. \\
Empresa de transporte fluvial & Responsabilidad solidaria como dueña de la cosa riesgosa (lancha) y empleadora de los responsables directos que actuaron en ejercicio de su actividad dependiente. \\
Compañía de seguros & Legitimada pasiva por el contrato de seguro que la vincula con la empresa, en virtud del Art. 118 de la Ley de Seguros. \\
\bottomrule
\end{longtable}

\section{Competencia}

En el caso analizado se optó por la Justicia Nacional, con fundamento en el domicilio legal de la empresa en la Ciudad de Buenos Aires. Existían, sin embargo, cuatro foros posibles:

\begin{itemize}
    \item Tribunales nacionales ordinarios (CABA).
    \item Tribunales federales.
    \item Tribunales de San Isidro (Provincia de Buenos Aires).
    \item Tribunales arbitrales.
\end{itemize}

La opción por la competencia federal existía porque el caso involucra un abordaje ocurrido en un río, materia prevista en la Ley de Navegación. Cabe señalar que, en el caso concreto, tramitaba en paralelo un juicio contra la misma empresa en sede federal en San Isidro, mientras que el caso bajo análisis se llevaba en justicia nacional, sin que los jueces intervinientes lograran ponerse de acuerdo respecto del conflicto de competencia planteado.

\subsection{Competencia penal: no es prorrogable}

\begin{important}
El hecho penal ocurrió en San Isidro, y esa competencia no puede prorrogarse, modificarse ni elegirse: la competencia penal se determina exclusivamente por el lugar en que se produjo el hecho.
\end{important}

\subsection{Instrucción penal preparatoria en la Provincia de Buenos Aires}

En la Provincia de Buenos Aires, a diferencia de la Ciudad de Buenos Aires, la investigación penal preparatoria está a cargo del fiscal, quien se ocupa de asentar los hechos y reunir las pruebas necesarias para determinar si existió un hecho delictivo y si este puede atribuirse a una persona determinada. No obstante que la instrucción es dirigida por el fiscal, existe siempre un juez de garantías, cuya función consiste en velar por que la instrucción respete las garantías procesales de las partes. Cuando la parte querellante o la defensa consideran que alguna actuación no se ajustó a derecho, pueden recurrir ante dicho juez de garantías.

\chapter{Citación en garantía del asegurador (Art. 118 Ley de Seguros)}

\section{Concepto}

\begin{definition}{Citación en garantía}
El asegurador no es responsable directo del hecho dañoso. Sin embargo, para el supuesto de que los demandados resulten condenados y no abonen la indemnización, el seguro constituye la garantía de pago. La citación en garantía es, entonces, el acto procesal mediante el cual se incorpora al asegurador como parte del proceso, a fin de que responda por la condena en la medida del seguro contratado.
\end{definition}

\section{Origen de la obligación del seguro}

La obligación del asegurador surge del contrato de seguro. Existen dos nociones relevantes: por un lado, determinadas actividades imponen legalmente la obligación de contar con seguro, como ocurre con el transporte público de pasajeros; por otro lado, la circulación de vehículos en general tampoco puede realizarse sin contar con la cobertura correspondiente. A través del contrato de seguro se contrata la cobertura de un riesgo determinado; cuando ese riesgo se concreta, se produce el siniestro.

\section{Necesidad de citar al asegurador desde el inicio}

\begin{important}
Para que la eventual condena alcance al asegurador, este debe haber sido parte en el proceso. Si el asegurador no fue citado como parte, la sentencia no podrá alcanzarlo ni podrá ser condenado a pagar. Por ello, la primera resolución que dicta el juez al pronunciarse sobre la demanda —``téngase por presentado y por parte''— resulta determinante, ya que fija a quiénes alcanzará la sentencia.
\end{important}

Del mismo modo en que, si las partes demandadas supieran que cuentan con cobertura de seguro, lo primero que harían sería denunciar el hecho a su aseguradora para que esta se presente a defenderse, la parte actora tiene interés en que el asegurador sea citado desde el inicio del proceso, precisamente porque, para que la condena lo alcance, debe haber tenido la posibilidad de ser parte y de ejercer su defensa.

\section{Texto normativo aplicable}

\begin{formula}{Art. 118 --- Ley de Seguros N.\textsuperscript{o} 17.418 --- Citación en garantía}
\begin{itemize}
    \item El damnificado puede citar en garantía al asegurador hasta que se reciba la causa a prueba.
    \item En tal caso debe interponer la demanda ante el juez del lugar del hecho o del domicilio del asegurador.
    \item El asegurador citado en garantía puede oponer a la víctima las defensas nacidas con anterioridad al siniestro.
    \item Si el asegurador no comparece, es representado por el asegurado.
    \item La sentencia dictada contra el asegurado alcanza al asegurador en la medida del seguro.
\end{itemize}
\end{formula}

\chapter{Los hechos del caso}

\section{Relato del accidente}

El relato de los hechos incorporado a la demanda reproduce la transcripción de las actuaciones del expediente penal correccional de San Isidro.

\begin{example}{Relato de los hechos (versión del expediente penal)}
``En el día X, aproximadamente a las [hora], la señora X y yo nos encontrábamos navegando en un kayak por el Río Tigre a la altura de la estación fluvial, habiendo partido de las instalaciones del club C. En esa circunstancia, la lancha L, que se encontraba amarrada como último de un amarre de varias embarcaciones, inició una maniobra de retroceso con la finalidad de realizar su zarpada, invirtiendo nuestra embarcación. A causa de la colisión, el kayak en el que nos encontrábamos dio vuelta campana y la señora X fue succionada por la turbulencia generada por las hélices de la lancha, produciéndole gravísimas lesiones que provocaron su muerte días después.''

``Según el relato de testigos presentes en el lugar, recogido por la actuación del personal de la Prefectura Naval Argentina y un video grabado desde el celular de uno de ellos, la maniobra por la lancha se realizó en violación al procedimiento usual, toda vez que el marinero se encontraba charlando con el patrón en lugar de estar apostado en la popa, verificando las condiciones adecuadas y dando las indicaciones para la zarpada más segura.''

``Al momento del siniestro contábamos con todos los elementos reglamentarios de seguridad para navegar y lo hacíamos en la zona habilitada para las embarcaciones a remo, observando las reglas de tránsito fluvial, lo que está acreditado por los testigos y el personal del club que autorizó nuestra salida.''
\end{example}

% TODO: el apunte original no consigna la fecha ni la hora exactas del hecho, ni los nombres de la víctima, del club ni de la lancha, reemplazados por placeholders (X, C, L, [hora]).

El último párrafo del relato reviste especial importancia, en tanto anticipa la defensa que previsiblemente opondrían los demandados (sosteniendo que el hecho no les fue imputable y que las víctimas se colocaron en una situación de riesgo) y acredita, en cambio, que los actores se encontraban en zona habilitada, eran socios del club con habilitación para navegar a remo y cumplían con las normas de tránsito fluvial vigentes. Se señala, asimismo, que el puerto fluvial constituye una zona de tránsito compartido por embarcaciones de distintos clubes que deben necesariamente cruzarlo, y que no existe, en los hechos, un representante de la Prefectura que ordene efectivamente el tránsito fluvial.

\section{Hechos del accidente vs. hechos que fundan los rubros}

A los fines de la redacción de la demanda, corresponde distinguir dos categorías de hechos que deben narrarse:

\begin{itemize}
    \item \textbf{Hechos del accidente}: la mecánica del siniestro, su desarrollo y los sujetos que intervinieron.
    \item \textbf{Hechos que fundan los rubros indemnizatorios}: las circunstancias personales de la víctima —su profesión, sus ingresos, su familia, sus actividades— que permiten justificar cada uno de los rubros reclamados.
\end{itemize}

\begin{note}
La existencia de hijos, el tiempo de convivencia de la pareja, el sufrimiento padecido, los deportes u otras actividades que practicaba la víctima, o el modo en que se organizaba su vida compartida, constituyen también cuestiones de hecho. Es una decisión estratégica del profesional interviniente ubicar esos hechos en el capítulo de hechos del accidente o distribuirlos en cada uno de los rubros indemnizatorios que se reclaman.
\end{note}

\chapter{Rubros indemnizatorios}

\section{Tipos de daño}

El daño susceptible de reclamo puede clasificarse en dos grandes categorías: daño patrimonial y daño extrapatrimonial.

\begin{longtable}{>{\raggedright\arraybackslash}p{0.45\textwidth} >{\raggedright\arraybackslash}p{0.45\textwidth}}
\toprule
\textbf{Daño patrimonial} & \textbf{Daño extrapatrimonial} \\
\midrule
\endfirsthead
\toprule
\textbf{Daño patrimonial} & \textbf{Daño extrapatrimonial} \\
\midrule
\endhead
Daño emergente: kayak destruido. & Daño moral: angustia, pena, tristeza permanente. \\
Lucro cesante: ingresos que la víctima dejó de percibir. & Daño psicológico: perturbación de la psiquis, tratable y cuantificable. \\
Pérdida de chance: expectativas futuras no concretadas. & Daño estético (si correspondiera). \\
Gastos médicos, de internación, velatorio y sepelio. & --- \\
\bottomrule
\end{longtable}

\section{Daño moral}

\begin{definition}{Daño moral}
El daño moral es un daño inmaterial: una tristeza, una angustia o una pena que no desaparece. En algunos supuestos requiere una comprobación mayor; en otros —como la pérdida de un ser querido— no requiere mayor comprobación, en la medida en que se entiende que toda persona sufre la pérdida de un familiar. Dado que la reparación integral, entendida como la vuelta de las cosas al estado anterior, resulta imposible en estos casos, el derecho ha encontrado en el dinero el único medio disponible de reparación.
\end{definition}

\section{Daño psicológico vs. daño moral}

\begin{definition}{Daño psicológico}
El daño psicológico consiste en la perturbación producida en la psiquis de la persona, y es tratable. Un perito puede establecer, por ejemplo, que la reparación de ese daño requiere una terapia de determinada frecuencia semanal durante un período de tiempo, y a partir del valor de las sesiones se obtiene el monto correspondiente. El daño moral no se confunde con este rubro, en tanto no admite ese tipo de reparación específica, al no existir manera de precisarlo mediante un tratamiento determinado.
\end{definition}

\begin{longtable}{>{\raggedright\arraybackslash}p{0.45\textwidth} >{\raggedright\arraybackslash}p{0.45\textwidth}}
\toprule
\textbf{Daño moral} & \textbf{Daño psicológico} \\
\midrule
\endfirsthead
\toprule
\textbf{Daño moral} & \textbf{Daño psicológico} \\
\midrule
\endhead
Daño inmaterial, espiritual, afectivo. & Perturbación de la psiquis de la persona. \\
Angustia, pena, tristeza permanentes. & Es tratable con terapia psicológica. \\
En casos de muerte de familiar: no requiere prueba pericial específica (se presume). & Requiere prueba pericial: el perito dictamina qué tratamiento necesita y por cuánto tiempo. \\
Difícil de cuantificar con exactitud. & Cuantificable: sesiones por valor por sesión por tiempo de tratamiento. \\
\bottomrule
\end{longtable}

\section{Lucro cesante vs. pérdida de chance}

El lucro cesante es lo que la víctima deja de percibir de manera concreta y mensurable, es decir, el ingreso mensual que efectivamente percibía. La pérdida de chance, en cambio, es aquello que la víctima eventualmente podría haber llegado a ganar, sin ser tan concreto, pero sí susceptible de estimación.

\begin{example}{Estimación según la actividad de la víctima}
Si la víctima trabajaba en relación de dependencia y percibía un ingreso mensual determinado, la expectativa de cálculo se proyecta hasta la edad jubilatoria. Si, en cambio, desarrollaba una labor profesional en crecimiento, la expectativa debe guardar relación con ese crecimiento proyectado. Si tenía una actividad comercial en auge o era deportista, las circunstancias son distintas, y en todos los casos deben ser enunciadas y probadas.
\end{example}

\section{Estrategia probatoria en el daño moral: el componente emotivo}

\begin{note}
La determinación del daño moral involucra, además de la cuestión técnica, una dimensión estratégica vinculada con la manera en que se presenta a la víctima. Existe un componente emocional, ideológico y de prejuicios sociales con el que el profesional debe trabajar al construir su estrategia probatoria: mostrar la magnitud real de la pérdida para una persona determinada, considerando sus circunstancias particulares, incide en la valuación del daño.
\end{note}

A modo de ilustración: si el actor practicaba un deporte de manera amateur y como consecuencia del hecho perdió la movilidad de un miembro, resulta pertinente ofrecer testigos, acreditar su condición de socio del club deportivo desde hace años y exhibir constancias de su participación, a fin de demostrar que la pérdida de esa actividad representa, para esa persona en concreto, un daño moral mayor que el que representaría para una persona sedentaria, incidiendo ello en la valuación del rubro.

\section{La competencia y el componente emotivo del juez}

La elección del foro también puede analizarse desde la perspectiva del componente emotivo que el caso genera en el juez. En el caso analizado, se señaló que una persona ajena a las actividades náuticas y residente en la Ciudad de Buenos Aires no tiene el mismo vínculo con el río que quien reside en la zona del Tigre, de modo que, en San Isidro, el hecho resulta más cercano y habitual para el juzgador, no porque este sea más sensible, sino porque el tema le resulta más próximo. Se trata de un componente emotivo que el profesional debe saber identificar y manejar al momento de definir su estrategia procesal.

\chapter{Relación entre el proceso civil y el proceso penal}

No resulta necesario aguardar la resolución del proceso penal para promover la demanda civil, en tanto el daño se produjo con independencia de aquel. Lo que aporta la resolución penal es que los hechos allí acreditados quedan, en los hechos, prácticamente probados para el proceso civil. Si, por el contrario, el hecho no resultara calificado como delito, ello incide en un sentido diferente, pero el daño sufrido por la víctima no se modifica por esa circunstancia. Que los demandados resulten condenados en sede penal favorece la posición de la parte actora en sede civil; en el caso analizado, el patrón y el marinero fueron efectivamente condenados en sede penal, circunstancia que no se verificó respecto de la empresa.

\chapter{Medios de prueba}

\section{Momento de ofrecimiento}

\begin{important}
En el proceso civil y comercial, la prueba debe ofrecerse con la demanda y con la contestación. Si no se ofrece en ese momento, no existe posibilidad de hacerlo con posterioridad. Este régimen difiere del previsto en otros fueros, como el derecho laboral, donde ciertas pruebas —con excepción de la documental— pueden ofrecerse en un momento distinto.
\end{important}

\section{Prueba documental}

\begin{definition}{Prueba documental}
La prueba documental es todo instrumento o soporte que acredita o demuestra la existencia de un hecho, de una relación jurídica o de una circunstancia determinada. No se limita a documentos en soporte papel, sino que comprende cualquier tipo de soporte que contenga información relevante.
\end{definition}

Como ejemplos de prueba documental mencionados en la clase se señalaron: el acta de matrimonio, el acta constitutiva de la empresa, el permiso de navegación, las filmaciones realizadas desde un teléfono celular y el acta de cierre de mediación.

\begin{important}
El instrumento que no fue acompañado junto con la demanda no puede presentarse con posterioridad. No existe posibilidad de subsanar esa omisión en una etapa procesal posterior, salvo que se configure la figura del hecho nuevo.
\end{important}

\section{Prueba testimonial, pericial e informativa}

\begin{note}
Estas tres modalidades probatorias serán objeto de una clase especial, en la cual se practicará la formulación de preguntas, se determinará quién las realiza y de qué manera, mediante una actuación simulada.
\end{note}

\chapter{Preclusión procesal y hecho nuevo}

\section{Concepto de preclusión}

\begin{definition}{Preclusión procesal}
La preclusión es el cierre de una etapa procesal. El proceso puede representarse como una cadena de eslabones, en la que cada acto procesal constituye un eslabón que debe cerrarse para poder unirse al siguiente y continuar así la cadena que se extiende desde el inicio del proceso hasta el dictado de la sentencia.
\end{definition}

\section{Cierre de cada acto procesal}

Cada acto procesal se cierra de dos maneras posibles: cuando la parte cumple el acto, o cuando vence el plazo previsto para cumplirlo.

\begin{example}{El caso trampa de la preclusión}
Si una parte cuenta con quince días para contestar una demanda y opta por hacerlo el día tres, ya no podrá completar ni modificar esa contestación en los días restantes, aun cuando el plazo original no hubiera vencido. Si, por ejemplo, omitió ofrecer prueba al contestar en el día tres, no podrá subsanar esa omisión, pues el acto ya se cerró: al haberlo ejercido anticipadamente, la parte perdió la posibilidad de utilizar los días restantes del plazo.
\end{example}

\section{La única excepción: ampliación de la demanda}

\begin{important}
La ampliación de la demanda constituye el único acto procesal que no precluye por su sola realización. Al presentar la demanda, esta todavía no ha sido notificada a la contraparte, por lo que se considera un acto no concluido, propio exclusivamente de la parte actora, quien conserva la posibilidad de ampliarlo. Una vez que la demanda es notificada a la contraparte, se traba la litis y precluye definitivamente la posibilidad de ampliarla.
\end{important}

\section{Hecho nuevo}

\begin{formula}{Hecho nuevo: requisitos y plazos}
El hecho nuevo constituye la excepción a la regla según la cual toda la prueba debe ofrecerse junto con la demanda, y permite incorporar con posterioridad una prueba que antes no se conocía. Para que un hecho pueda calificarse como hecho nuevo deben reunirse los siguientes requisitos:
\begin{itemize}
    \item Que no se conociera al momento de presentar la demanda, o que se hubiera producido con posterioridad a esta.
    \item Que tenga relevancia para la resolución del caso.
\end{itemize}
En cuanto a los plazos:
\begin{itemize}
    \item Plazo para presentarlo: cinco días desde que se tomó conocimiento del hecho.
    \item Límite máximo: hasta la audiencia prevista en el artículo 360 del CPCCN, o hasta cinco días antes de su celebración.
    \item Si el hecho aparece con posterioridad a dicha audiencia, su tratamiento se canaliza mediante el sistema de recursos, cuestión que excede el alcance de esta clase.
\end{itemize}
\end{formula}

\chapter{Cierre de la clase y próximos pasos}

Se estableció que en las clases siguientes se profundizará en los siguientes contenidos:

\begin{itemize}
    \item Estudio de las excepciones procesales (defensas de forma y de fondo).
    \item Prueba testimonial, informativa y pericial, con actuación práctica.
    \item Elaboración, por parte de cada alumno, de su propio modelo de demanda, para su corrección y exposición oral.
\end{itemize}

\chapter{Notas Cornell: síntesis de la clase}

El siguiente cuadro reúne, mediante el método Cornell, las preguntas centrales que permiten repasar los conceptos desarrollados a lo largo de la clase, junto con sus respuestas y una síntesis final.

\begin{table}[H]
\centering
\begin{tabularx}{\textwidth}{
    >{\raggedright\arraybackslash}p{0.28\textwidth}
    >{\raggedright\arraybackslash}X
}
\toprule
\textbf{Preguntas / palabras clave} &
\textbf{Notas / respuestas} \\
\midrule

\textbf{¿Qué es la citación en garantía?} &
Es el acto procesal por el cual el actor incorpora al asegurador del demandado como parte en el juicio, en virtud del Art. 118 de la Ley de Seguros N.\textsuperscript{o} 17.418. Sirve para que la condena eventualmente dictada contra el asegurado sea oponible y ejecutable también contra la aseguradora. Si el asegurador no fue citado como parte, no puede ser condenado ni obligado a pagar. \\

\textbf{¿Qué dice el Art. 330 CPCCN?} &
Establece los requisitos formales de la demanda: nombre y domicilio del demandante y del demandado, la cosa demandada con exactitud, los hechos explicados claramente, el derecho sucintamente expuesto y la petición en términos claros y positivos, con indicación del monto reclamado o justificación de su imposibilidad. \\

\textbf{¿Cuál es la diferencia entre carga procesal y obligación?} &
La obligación es el deber jurídico de hacer o no hacer, exigible judicialmente y coercible. La carga procesal no tiene sanción externa: si no se ejerce, el perjudicado es el propio sujeto, que pierde la oportunidad de hacerlo. En el proceso civil todas son cargas: ofrecer prueba, plantear excepciones, contestar traslados. No ejercerlas no es un delito, pero reduce las probabilidades de obtener una sentencia favorable. \\

\textbf{¿Qué es la legitimación activa y la pasiva?} &
La legitimación activa es la calidad que habilita a una persona a estar en juicio como actor (quién puede demandar y por qué). La legitimación pasiva es la calidad que habilita a una persona a ser demandada (quiénes son responsables y por qué título jurídico). Plantear expresamente ambas desde el inicio evita o debilita la excepción de falta de legitimación que intentará la contraparte. \\

\textbf{¿Qué diferencia hay entre daño moral y daño psicológico?} &
El daño moral es inmaterial: angustia, pena, tristeza permanente. En caso de muerte de familiar próximo no requiere prueba pericial específica (se presume). El daño psicológico es la perturbación de la psiquis, tratable con terapia y cuantificable mediante peritaje: el perito indica el tipo y duración del tratamiento y se multiplica por el costo de cada sesión para obtener el monto. \\

\textbf{¿Qué es la preclusión procesal y cuándo opera?} &
La preclusión es el cierre de una etapa procesal. Opera de dos maneras: (1) cuando el sujeto cumple el acto antes de vencer el plazo (el acto se clausura al realizarse, aunque queden días disponibles); y (2) cuando vence el plazo sin haber cumplido el acto. La única excepción es la ampliación de la demanda, que no precluye hasta que la demanda sea notificada a la contraparte y se trabe la litis. \\

\textbf{¿Qué es un hecho nuevo y cuáles son sus requisitos?} &
Es un hecho o prueba sobreviniente que puede incorporarse al proceso después de la demanda. Requisitos: (1) que no se conociera al momento de presentar la demanda, o (2) que se haya producido con posterioridad, y (3) que sea relevante para la resolución del caso. Plazo: cinco días desde que se conoce el hecho. Límite: hasta la audiencia del Art. 360 CPCCN. \\

\textbf{¿Cómo se diferencia lucro cesante de pérdida de chance?} &
El lucro cesante es lo que se dejó de percibir de manera concreta y mensurable (los ingresos mensuales que la víctima ganaba). La pérdida de chance es la expectativa futura de ganancia que no es concreta pero sí estimable (lo que la víctima podría haber llegado a ganar según su trayectoria profesional, crecimiento laboral, etc.). Ambos son rubros patrimoniales que deben enunciarse y probarse. \\

\textbf{¿Cuándo deben ofrecerse las pruebas en el proceso civil?} &
En el proceso civil y comercial (CPCCN), la prueba documental debe acompañarse con la demanda o la contestación sin excepción. Las demás pruebas (testimonial, pericial, informativa) también deben ofrecerse con esos escritos. Si no se ofrecen en ese momento, precluye la oportunidad de hacerlo y no hay posibilidad de incorporarlas posteriormente, salvo la figura del hecho nuevo. \\

\midrule

\multicolumn{2}{p{\textwidth}}{
\textbf{Resumen:}
Esta clase expuso en forma práctica los pilares del proceso civil: cómo se redacta y estructura una demanda conforme al Art. 330 CPCCN; cómo se determina y se justifica la competencia; cómo se identifica a cada demandado según su título de responsabilidad —directa, solidaria o contractual—; y cómo se incorpora al asegurador mediante la citación en garantía del Art. 118 de la Ley de Seguros, herramienta indispensable en todo juicio por daños. A continuación, se desarrollaron los rubros indemnizatorios: la distinción entre daño patrimonial (daño emergente, lucro cesante, pérdida de chance, gastos médicos y funerarios) y daño extrapatrimonial (daño moral y daño psicológico), con sus diferencias en materia de prueba y cuantificación. Se destacó que el trabajo del abogado no es solo técnico sino también estratégico y comunicacional: saber a qué juez le hablo, qué prueba construyo y cómo presento a mi cliente como la víctima son parte del oficio. Finalmente, se analizaron los conceptos procesales de preclusión —cómo el ejercicio prematuro de un acto clausura la etapa, sin importar el plazo residual— y hecho nuevo —la válvula de escape para incorporar prueba sobreviniente—, cerrando así el ciclo completo de la acción judicial desde su presentación hasta la etapa probatoria.
} \\

\bottomrule
\end{tabularx}
\end{table}

\end{document}
